\documentclass[11pt,a4paper]{article}
\usepackage[margin=1in]{geometry}
\usepackage[T1]{fontenc}
\usepackage[utf8]{inputenc}
\usepackage{lmodern}
\usepackage{microtype}
\usepackage{amsmath,amssymb,mathtools}
\usepackage[hidelinks]{hyperref}
\usepackage{fancyhdr}
\setlength{\headheight}{14pt}
\pagestyle{fancy}
\fancyhf{}
\lhead{Ernst Steinitz}
\rhead{English}
\cfoot{\thepage}
\setlength{\parindent}{1.35em}
\setlength{\parskip}{0.35em}
\allowdisplaybreaks
\newcommand{\K}{\mathfrak K}
\newcommand{\Kl}{\operatorname{Kl}}
\newcommand{\Zkl}{\operatorname{Zkl}}
\newcommand{\Kkl}{\operatorname{Kkl}}
\newcommand{\Rg}{\operatorname{Rg}}
\begin{document}

\section*{Rectangular Systems and Modules in Algebraic Number Fields. I}
\begin{center}
By \emph{Ernst Steinitz} in Breslau.
\end{center}

\subsection*{Introduction}

\textbf{1.} In what follows we fix an arbitrary finite algebraic number field \(\K\). Unless otherwise stated, only numbers and ideals of \(\K\) are considered, and they are integral numbers and integral ideals.

\textbf{2.} Let
\[
 A=(a_{ik}),\qquad B=(b_{kl})
\]
be two rectangular systems of integral numbers from \(\K\). If the number of columns of \(A\) is equal to the number of rows of \(B\), their product is the rectangle \(C=(c_{il})\) defined by
\[
 c_{il}=\sum_k a_{ik}b_{kl}.
\]
Then \(C\) has as many rows as \(A\) and as many columns as \(B\). This multiplication is associative, but not necessarily commutative. If two rectangles \(A,B\) satisfy a relation
\[
 B=PAQ,
\]
then \(B\) is called divisible by \(A\). If, in addition, \(A\) is divisible by \(B\), then \(A\) and \(B\) are equivalent.

\textbf{3.} The first question treated here is the criterion for equivalence of two systems \(A,B\); in a second article the more general question of divisibility of \(B\) by \(A\) is to be treated. Over the rational field the answer is known. Over an arbitrary algebraic number field the final results remain the same for square systems of nonzero determinant, although the proof requires other tools. Once one allows nonsquare systems, or square systems of determinant \(0\), the results themselves change.

\textbf{4.} Left and right divisibility are special cases. \(B\) is left-divisible by \(A\) if \(B=PA\), and right-divisible by \(A\) if \(B=AQ\). Correspondingly, two systems are left- or right-equivalent if each is left- or right-divisible by the other.

\subsection*{Preliminaries}

\textbf{5.} We recall the ideal-theoretic facts and notation that will be used. The nonzero ideals of \(\K\) are divided into finitely many classes; let their number be \(h\). The class of the product of two ideals \(\mathfrak a,\mathfrak b\) depends only on the classes of the factors:
\[
 \Kl(\mathfrak a\mathfrak b)=\Kl(\mathfrak a)\Kl(\mathfrak b). \tag{1}
\]
The \(h\) classes form an abelian group. The identity is the principal class, consisting of principal ideals. Fractional ideals are assigned to the same classes in such a way that (1) remains valid.

\textbf{6.} Any finite family of integral ideals \(\mathfrak a_1,\ldots,\mathfrak a_n\) has a greatest common divisor
\[
 (\mathfrak a_1,\ldots,\mathfrak a_n)=\mathfrak t,
\]
which is a multiple of every common divisor. If all ideals are zero, then \(\mathfrak t=0\). For numbers \(a_1,\ldots,a_n\) this means the gcd of the principal ideals generated by them; the expressions
\[
 a_1x_1+\cdots+a_nx_n
\]
represent precisely the numbers divisible by \((a_1,\ldots,a_n)\). For integral or fractional ideals,
\[
 (\mathfrak a_1,\ldots,\mathfrak a_n)\mathfrak q
   =(\mathfrak a_1\mathfrak q,\ldots,\mathfrak a_n\mathfrak q) \tag{2}
\]
and more generally
\[
 (\mathfrak a_1,\ldots,\mathfrak a_m)(\mathfrak b_1,\ldots,\mathfrak b_n)
 =(\mathfrak a_i\mathfrak b_j\mid 1\leq i\leq m,\;1\leq j\leq n). \tag{3}
\]
Even an infinite system of integral ideals has a greatest common divisor, already determined by some finite subsystem.

\textbf{7.} Let \(\mathfrak p_1,\ldots,\mathfrak p_l\) be distinct prime ideals and \(e_1,\ldots,e_l\) nonnegative rational integers. In every ideal class there are ideals that contain each \(\mathfrak p_i\) exactly to the prescribed exponent \(e_i\). In particular, in every class one can find ideals divisible by a given ideal, and also ideals relatively prime to a given ideal.

\textbf{8.} Let \(a_1,\ldots,a_n\) be not all zero and put
\[
 (a_1,\ldots,a_n)=\mathfrak t.
\]
As \(c\) ranges through the nonzero elements of \(\K\), the systems \(a_1c,\ldots,a_nc\) are precisely the systems proportional to \(a_1,\ldots,a_n\), and
\[
 (a_1c,\ldots,a_nc)=\mathfrak t(c)
\]
ranges over the ideals in the class \(\Kl(\mathfrak t)\). Thus every system is proportional to an integral system whose elements have no nontrivial common divisor.

\subsection*{Elementary Divisors, Row Class and Column Class}

\textbf{9.} If \(A\) is a rectangular system of integral numbers from \(\K\), the greatest common divisor \(\mathfrak d_k\) of its \(k\)-rowed minors is called the \(k\)-th determinant divisor of \(A\). If \(k\) is larger than the number of rows or columns, set \(\mathfrak d_k=0\). Laplace's determinant theorem implies that \(\mathfrak d_k\) divides \(\mathfrak d_{k+1}\). If \(\mathfrak d_{r+1}\) is the first zero determinant divisor, \(r\) is the rank of \(A\):
\[
 r=\Rg A.
\]
With \(\mathfrak d_0=(1)\), the ideals
\[
 \mathfrak e_k=\mathfrak d_k\mathfrak d_{k-1}^{-1}
\]
are the elementary divisors of \(A\).

\textbf{10.} If
\[
 B=PAQ,
\]
then each determinant divisor of \(B\) is divisible by the corresponding determinant divisor of \(A\). Consequently equivalent rectangular systems have the same determinant divisors, hence the same elementary divisors. This is the first condition for equivalence.

\textbf{11.} Let \(A=(a_{ik})\) be an integral \(m\times n\) rectangle of rank \(r\geq1\). Form a new rectangle \(A'\) from all \(r\)-minors of \(A\), arranged so that its rows correspond to choices of \(r\) rows and its columns to choices of \(r\) columns. This rectangle has rank \(1\). The gcds of the nonzero rows of \(A'\) therefore belong to one ideal class; this is the row class
\[
 \Zkl A.
\]
The gcds of the columns similarly determine the column class
\[
 \Kkl A.
\]

\textbf{12.} If \(\mathfrak d_r\) is the \(r\)-th determinant divisor of a rectangle \(A\) of rank \(r\), then
\[
 \Zkl A\cdot \Kkl A=\Kl(\mathfrak d_r). \tag{4}
\]
Indeed, if one nonzero element of the rectangle of \(r\)-minors \(A'\) is written as \(uv\), then every row of \(A'\) is proportional to the same system of \(v\)'s and every column to the same system of \(u\)'s. The common divisor of all elements of \(A'\) is exactly \(\mathfrak d_r\).

\textbf{13.} If \(B=PA\) and \(A,B\) have the same rank, then they have the same row class:
\[
 \Zkl B=\Zkl A.
\]
Likewise, if \(B=AQ\) and the ranks agree, the column class is preserved. This follows by expressing the \(r\)-minors of \(B\) as linear combinations of the \(r\)-minors of \(A\) by the general multiplication theorem for determinants.

\textbf{14.} If \(A\) and \(B=PAQ\) are equivalent rectangular systems of rank \(r\geq1\), then they have the same determinant divisors, the same elementary divisors, and the same row and column classes:
\[
 \Zkl A=\Zkl B,\qquad \Kkl A=\Kkl B. \tag{5}
\]
This is the second condition for equivalence.

\subsection*{Linear Forms, Equations and Congruences}

\textbf{15.} Consider \(m\) linear forms in \(n\) unknowns,
\[
 y_i=a_{i1}x_1+\cdots+a_{in}x_n\qquad(i=1,\ldots,m), \tag{6}
\]
with integral coefficients in \(\K\). Let \(\mathfrak t\) be the gcd of all coefficients \(a_{ik}\). If all \(x_k\) are chosen divisible by an ideal \(\mathfrak u\), then every \(y_i\) is divisible by \(\mathfrak t\mathfrak u\). If the coefficient system \(A=(a_{ik})\) has rank at least \(2\), the \(x_k\) can be chosen so that
\[
 (y_1,\ldots,y_m)=\mathfrak t\mathfrak u. \tag{7}
\]
The proof uses a finite set of prime ideals \(\mathfrak p_1,\ldots,\mathfrak p_l\) containing the prime factors of \(\mathfrak t\mathfrak u\) and of one nonzero two-rowed determinant. First the \(x_k\) are chosen by simultaneous congruences
\[
 x_k\equiv \xi_{ik}\pmod{\mathfrak p_i^{e_i+f_i+1}}
\]
so that at each \(\mathfrak p_i\) the gcd of the values has the prescribed exponent. Then a nonzero determinant
\[
 \Delta=a_{11}a_{22}-a_{12}a_{21}
\]
is used to alter two variables without destroying those congruences, eliminating all remaining prime factors. This proves (7). For \(\mathfrak u=(1)\), the values of the linear forms can therefore be made to have the same gcd as the coefficients.

\textbf{16.} Let \(A\) be a rectangle with \(n\) columns and \(m\leq n-2\) rows, of rank \(m\), and let \(\mathfrak b\) be its \(m\)-th determinant divisor. If one adjoins a new row \(x_1,\ldots,x_n\), with all entries divisible by a prescribed ideal \(\mathfrak u\), the row can be chosen so that the \((m+1)\)-st determinant divisor of the enlarged rectangle is
\[
 \mathfrak b\mathfrak u.
\]
Without a divisibility condition on the new row, the determinant divisor can be kept equal to \(\mathfrak b\). Iterating gives: a rectangle with \(n\) columns and \(m<n\) rows whose \(m\)-minors are coprime can be completed by further rows to a square system of determinant \(1\).

\textbf{17.} If the coefficient system in (6) has rank at most \(n-2\), the equations
\[
 y_i=0\qquad(i=1,\ldots,m)
\]
have a solution in coprime numbers. Choose two linearly independent integral solutions \(\alpha=(\alpha_k)\) and \(\beta=(\beta_k)\), arrange that their gcds are relatively prime, and then choose \(\lambda,\mu\) so that the numbers
\[
 \alpha_k\lambda+\beta_k\mu
\]
are coprime.

\textbf{18.} If the coefficient system in (6) has rank at most \(n-1\), and if \(\mathfrak m\ne0\) is an ideal, the congruences
\[
 y_i\equiv0\pmod{\mathfrak m}\qquad(i=1,\ldots,m) \tag{8}
\]
can be solved in coprime numbers. One first chooses a nontrivial integral solution of \(y_i=0\) whose gcd is relatively prime to \(\mathfrak m\), and then replaces its entries by congruent integral numbers modulo \(\mathfrak m\) so as to obtain an actually coprime solution of the congruences.

\textbf{19.} Let \(\mathfrak e_n\) be the \(n\)-th elementary divisor of the coefficient system \(A=(a_{ik})\) in (6). Every solution of the congruences (8) consists of numbers divisible by
\[
 \frac{\mathfrak m}{(\mathfrak m,\mathfrak e_n)}.
\]
If \(\mathfrak e_n=0\), this is trivial. Otherwise let \(\mathfrak b_n\) and \(\mathfrak b_{n-1}\) be the \(n\)-th and \((n-1)\)-st determinant divisors, so that
\[
 \mathfrak e_n=\mathfrak b_n\mathfrak b_{n-1}^{-1}.
\]
From any nonsingular subsystem of \(n\) equations, Cramer's rule expresses each \(x_k\) as a quotient whose denominator is an \(n\)-rowed determinant and whose numerator is divisible by \(\mathfrak b_{n-1}\mathfrak m\). Letting the nonzero \(n\)-determinants range over all such subsystems and using their gcd \(\mathfrak b_n\) yields the asserted divisibility.

\textbf{20.} The homogeneous linear congruences
\[
 \sum_{k=1}^{n}a_{ik}x_k\equiv0\pmod{\mathfrak m}\qquad(i=1,\ldots,m) \tag{9}
\]
are solvable in coprime numbers if and only if \(\mathfrak m\) divides the \(n\)-th elementary divisor \(\mathfrak e_n\) of the coefficient system. Necessity is Nr. 19. For sufficiency one reduces to \(\mathfrak m=\mathfrak e_n\), factors \(\mathfrak e_n\) into prime-ideal powers, solves the congruences for each prime-ideal power by Nr. 18, and combines the resulting solutions by simultaneous congruences. The final numbers may then be replaced by congruent integral numbers so that their gcd is \(1\).

\end{document}
